%%%%%%%%%%%%%%%%%%%%%%%%%%%%%%%%%%%%%%%%%%%%%%%%%%%%%%%%%%%%%%%%%%%%%%%%%%%%%%%%
% CHAPTER 8
%%%%%%%%%%%%%%%%%%%%%%%%%%%%%%%%%%%%%%%%%%%%%%%%%%%%%%%%%%%%%%%%%%%%%%%%%%%%%%%%

\chapter{Discrete-Time Fourier Transform and Frequency Analysis}

\begin{tcolorbox}[colback=blue!5!white,colframe=blue!70!black,title=Learning Objectives]

After studying this chapter, the reader will be able to

\begin{itemize}
\item Understand the need for DTFT.
\item Derive the DTFT and inverse DTFT.
\item Interpret the DTFT as a periodic spectrum.
\item Compute DTFTs of common sequences.
\item Understand magnitude and phase spectra.
\item Analyze frequency response of discrete-time systems.
\end{itemize}

\end{tcolorbox}

%%%%%%%%%%%%%%%%%%%%%%%%%%%%%%%%%%%%%%%%%%%%%%%%%%%%%%%%%%%%%%%%%%%%%%%%%%%%%%%%

\section{Introduction}

For continuous-time aperiodic signals, we use the CTFT. For discrete-time
aperiodic sequences, we use the \textbf{Discrete-Time Fourier Transform (DTFT)}.

The DTFT converts a sequence

\[
x[n]
\]

into a continuous function of frequency

\[
X(e^{j\omega}).
\]

%%%%%%%%%%%%%%%%%%%%%%%%%%%%%%%%%%%%%%%%%%%%%%%%%%%%%%%%%%%%%%%%%%%%%%%%%%%%%%%%

\section{DTFT Definition}

\[
\boxed{
X(e^{j\omega})=
\sum_{n=-\infty}^{\infty}
x[n]e^{-j\omega n}
}
\]

where

\[
-\pi\le\omega\le\pi.
\]

%%%%%%%%%%%%%%%%%%%%%%%%%%%%%%%%%%%%%%%%%%%%%%%%%%%%%%%%%%%%%%%%%%%%%%%%%%%%%%%%

\section{Inverse DTFT}

\[
\boxed{
x[n]=
\frac{1}{2\pi}
\int_{-\pi}^{\pi}
X(e^{j\omega})e^{j\omega n}\,d\omega
}
\]

The pair is written as

\[
\boxed{
x[n]\xleftrightarrow{\mathcal{F}}X(e^{j\omega})
}
\]

%%%%%%%%%%%%%%%%%%%%%%%%%%%%%%%%%%%%%%%%%%%%%%%%%%%%%%%%%%%%%%%%%%%%%%%%%%%%%%%%

\section{Why is the DTFT Periodic?}

Consider

\[
X(e^{j(\omega+2\pi)})=
\sum_{n=-\infty}^{\infty}
x[n]e^{-j(\omega+2\pi)n}.
\]

Since

\[
e^{-j2\pi n}=1
\]

for every integer \(n\),

\[
\boxed{
X(e^{j(\omega+2\pi)})=X(e^{j\omega})
}
\]

Hence the DTFT is periodic with period

\[
\boxed{
2\pi
}
\]

%%%%%%%%%%%%%%%%%%%%%%%%%%%%%%%%%%%%%%%%%%%%%%%%%%%%%%%%%%%%%%%%%%%%%%%%%%%%%%%%

\section{Frequency Axis}

The normalized frequency is

\[
\omega \text{ rad/sample}.
\]

One complete period is

\[
-\pi\le\omega\le\pi.
\]

The corresponding frequency in hertz is

\[
\boxed{
f=\frac{\omega}{2\pi}f_s
}
\]

where \(f_s\) is the sampling frequency.

%%%%%%%%%%%%%%%%%%%%%%%%%%%%%%%%%%%%%%%%%%%%%%%%%%%%%%%%%%%%%%%%%%%%%%%%%%%%%%%%

\section{DTFT of the Unit Sample}

Let

\[
x[n]=\delta[n].
\]

Then

\[
X(e^{j\omega})=
\sum_{n=-\infty}^{\infty}
\delta[n]e^{-j\omega n}.
\]

Only the term \(n=0\) survives:

\[
\boxed{
X(e^{j\omega})=1
}
\]

Thus,

\[
\boxed{
\delta[n]\xleftrightarrow{\mathcal{F}}1
}
\]

%%%%%%%%%%%%%%%%%%%%%%%%%%%%%%%%%%%%%%%%%%%%%%%%%%%%%%%%%%%%%%%%%%%%%%%%%%%%%%%%

\section{DTFT of a Shifted Impulse}

For

\[
x[n]=\delta[n-n_0],
\]

we obtain

\[
X(e^{j\omega})=
e^{-j\omega n_0}.
\]

Hence,

\[
\boxed{
\delta[n-n_0]
\xleftrightarrow{\mathcal{F}}
e^{-j\omega n_0}
}
\]

%%%%%%%%%%%%%%%%%%%%%%%%%%%%%%%%%%%%%%%%%%%%%%%%%%%%%%%%%%%%%%%%%%%%%%%%%%%%%%%%

\section{DTFT of the Unit Step}

The DTFT of \(u[n]\) is not absolutely convergent in the ordinary sense. It is
interpreted using generalized functions:

\[
\boxed{
u[n]
\xleftrightarrow{\mathcal{F}}
\pi\sum_{k=-\infty}^{\infty}\delta(\omega-2\pi k)
+\frac{1}{1-e^{-j\omega}}
}
\]

In practice, we usually work with the exponentially weighted sequence
\(a^n u[n]\).

%%%%%%%%%%%%%%%%%%%%%%%%%%%%%%%%%%%%%%%%%%%%%%%%%%%%%%%%%%%%%%%%%%%%%%%%%%%%%%%%

\section{DTFT of \(a^n u[n]\)}

Assume \(|a|<1\).

\[
X(e^{j\omega})=
\sum_{n=0}^{\infty}
a^n e^{-j\omega n}
\]

\[
=
\sum_{n=0}^{\infty}
(ae^{-j\omega})^n.
\]

This is a geometric series:

\[
\boxed{
X(e^{j\omega})=
\frac{1}{1-ae^{-j\omega}}
}
\]

with convergence condition

\[
|a|<1.
\]

%%%%%%%%%%%%%%%%%%%%%%%%%%%%%%%%%%%%%%%%%%%%%%%%%%%%%%%%%%%%%%%%%%%%%%%%%%%%%%%%

\section{Magnitude Spectrum}

For

\[
X(e^{j\omega})=
\frac{1}{1-ae^{-j\omega}},
\]

the magnitude is

\[
|X(e^{j\omega})|=
\frac{1}{|1-ae^{-j\omega}|}.
\]

Using

\[
|1-ae^{-j\omega}|^2=
1+a^2-2a\cos\omega,
\]

we get

\[
\boxed{
|X(e^{j\omega})|=
\frac{1}{\sqrt{1+a^2-2a\cos\omega}}
}
\]

%%%%%%%%%%%%%%%%%%%%%%%%%%%%%%%%%%%%%%%%%%%%%%%%%%%%%%%%%%%%%%%%%%%%%%%%%%%%%%%%

\section{Phase Spectrum}

\[
\boxed{
\angle X(e^{j\omega})=
-\tan^{-1}
\left(
\frac{a\sin\omega}{1-a\cos\omega}
\right)
}
\]

%%%%%%%%%%%%%%%%%%%%%%%%%%%%%%%%%%%%%%%%%%%%%%%%%%%%%%%%%%%%%%%%%%%%%%%%%%%%%%%%

\section{Rectangular Sequence}

Consider

\[
x[n]=
\begin{cases}
1, & 0\le n\le N-1\\
0, & \text{otherwise}
\end{cases}
\]

%%%%%%%%%%%%%%%%%%%%%%%%%%%%%%%%%%%%%%%%%%%%%%%%%%%%%%%%%%%%%%%%%%%%%%%%%%%%%%%%

\subsection{DTFT}

\[
X(e^{j\omega})=
\sum_{n=0}^{N-1}e^{-j\omega n}.
\]

Using the finite geometric-series formula,

\[
X(e^{j\omega})=
\frac{1-e^{-j\omega N}}{1-e^{-j\omega}}.
\]

Multiplying numerator and denominator by \(e^{-j\omega(N-1)/2}\),

\[
\boxed{
X(e^{j\omega})=
e^{-j\omega(N-1)/2}
\frac{\sin(N\omega/2)}{\sin(\omega/2)}
}
\]

%%%%%%%%%%%%%%%%%%%%%%%%%%%%%%%%%%%%%%%%%%%%%%%%%%%%%%%%%%%%%%%%%%%%%%%%%%%%%%%%

\section{Magnitude of Rectangular Sequence}

\[
\boxed{
|X(e^{j\omega})|=
\left|
\frac{\sin(N\omega/2)}{\sin(\omega/2)}
\right|
}
\]

This function is called the \textbf{Dirichlet kernel}.

%%%%%%%%%%%%%%%%%%%%%%%%%%%%%%%%%%%%%%%%%%%%%%%%%%%%%%%%%%%%%%%%%%%%%%%%%%%%%%%%

\section{Main-Lobe Width}

The first zeros occur at

\[
\omega=\pm\frac{2\pi}{N}.
\]

Hence the main-lobe width is

\[
\boxed{
\Delta\omega=\frac{4\pi}{N}
}
\]

A longer sequence produces a narrower main lobe and better frequency resolution.

%%%%%%%%%%%%%%%%%%%%%%%%%%%%%%%%%%%%%%%%%%%%%%%%%%%%%%%%%%%%%%%%%%%%%%%%%%%%%%%%

\section{DTFT of a Cosine}

Consider

\[
x[n]=\cos(\omega_0 n).
\]

Using Euler’s identity,

\[
\cos(\omega_0 n)=
\frac{1}{2}
\left(
e^{j\omega_0 n}+e^{-j\omega_0 n}
\right).
\]

The DTFT is

\[
\boxed{
X(e^{j\omega})=
\pi\sum_{k=-\infty}^{\infty}
\left[
\delta(\omega-\omega_0-2\pi k)
+\delta(\omega+\omega_0-2\pi k)
\right]
}
\]

Within \(-\pi\le\omega\le\pi\), there are two impulses at

\[
\boxed{
\omega=\pm\omega_0.
}
\]

%%%%%%%%%%%%%%%%%%%%%%%%%%%%%%%%%%%%%%%%%%%%%%%%%%%%%%%%%%%%%%%%%%%%%%%%%%%%%%%%

\section{Frequency Response of an LTI System}

For an LTI system,

\[
y[n]=x[n]*h[n].
\]

Taking the DTFT,

\[
\boxed{
Y(e^{j\omega})=
X(e^{j\omega})H(e^{j\omega})
}
\]

where

\[
\boxed{
H(e^{j\omega})=
\sum_{n=-\infty}^{\infty}
h[n]e^{-j\omega n}
}
\]

is the frequency response.

%%%%%%%%%%%%%%%%%%%%%%%%%%%%%%%%%%%%%%%%%%%%%%%%%%%%%%%%%%%%%%%%%%%%%%%%%%%%%%%%

\section{Example: Two-Point Averager}

Let

\[
h[n]=\delta[n]+\delta[n-1].
\]

Then

\[
H(e^{j\omega})=
1+e^{-j\omega}.
\]

Factorizing,

\[
H(e^{j\omega})=
2e^{-j\omega/2}\cos\frac{\omega}{2}.
\]

Therefore,

\[
\boxed{
|H(e^{j\omega})|=
2\left|\cos\frac{\omega}{2}\right|
}
\]

and

\[
\boxed{
\angle H(e^{j\omega})=-\frac{\omega}{2}.
}
\]

This is a low-pass filter.

%%%%%%%%%%%%%%%%%%%%%%%%%%%%%%%%%%%%%%%%%%%%%%%%%%%%%%%%%%%%%%%%%%%%%%%%%%%%%%%%

\section{Low-Pass Behavior}

Evaluate the magnitude:

At \(\omega=0\),

\[
|H(e^{j0})|=2.
\]

At \(\omega=\pi\),

\[
|H(e^{j\pi})|=0.
\]

Thus, low frequencies are passed and high frequencies are rejected.

%%%%%%%%%%%%%%%%%%%%%%%%%%%%%%%%%%%%%%%%%%%%%%%%%%%%%%%%%%%%%%%%%%%%%%%%%%%%%%%%

\section{High-Pass Example}

Let

\[
h[n]=\delta[n]-\delta[n-1].
\]

Then

\[
H(e^{j\omega})=
1-e^{-j\omega}
=
2je^{-j\omega/2}\sin\frac{\omega}{2}.
\]

Hence,

\[
\boxed{
|H(e^{j\omega})|=
2\left|\sin\frac{\omega}{2}\right|
}
\]

Now,

\[
|H(e^{j0})|=0,
\qquad
|H(e^{j\pi})|=2.
\]

This is a high-pass filter.

%%%%%%%%%%%%%%%%%%%%%%%%%%%%%%%%%%%%%%%%%%%%%%%%%%%%%%%%%%%%%%%%%%%%%%%%%%%%%%%%

\section{Bandwidth Interpretation}

For discrete-time systems,

\[
\omega=0
\]

represents DC, while

\[
\omega=\pi
\]

represents the highest representable frequency (Nyquist frequency).

The useful frequency range is

\[
\boxed{
0\le\omega\le\pi.
}
\]

%%%%%%%%%%%%%%%%%%%%%%%%%%%%%%%%%%%%%%%%%%%%%%%%%%%%%%%%%%%%%%%%%%%%%%%%%%%%%%%%

\section{Quick Check}

\begin{enumerate}
\item DTFT period = \(\boxed{2\pi}\).
\item \(a^n u[n]\leftrightarrow\boxed{\dfrac{1}{1-ae^{-j\omega}}}\).
\item Rectangular sequence DTFT = \(\boxed{\dfrac{\sin(N\omega/2)}{\sin(\omega/2)}}\).
\item Low-pass filter: \(\boxed{1+e^{-j\omega}}\).
\item High-pass filter: \(\boxed{1-e^{-j\omega}}\).
\end{enumerate}

%%%%%%%%%%%%%%%%%%%%%%%%%%%%%%%%%%%%%%%%%%%%%%%%%%%%%%%%%%%%%%%%%%%%%%%%%%%%%%%%

\begin{tcolorbox}[title=One-Minute Revision,colback=blue!5]

\[
X(e^{j\omega})=
\sum_{n=-\infty}^{\infty}
x[n]e^{-j\omega n}
\]

\[
x[n]=
\frac{1}{2\pi}
\int_{-\pi}^{\pi}
X(e^{j\omega})e^{j\omega n}d\omega
\]

\[
X(e^{j(\omega+2\pi)})=X(e^{j\omega})
\]

\[
a^n u[n]
\leftrightarrow
\frac{1}{1-ae^{-j\omega}}
\]

\[
|X(e^{j\omega})|=
\frac{1}{\sqrt{1+a^2-2a\cos\omega}}
\]

\[
H(e^{j\omega})=
\sum h[n]e^{-j\omega n}
\]

Low-pass:

\[
1+e^{-j\omega}
\]

High-pass:

\[
1-e^{-j\omega}
\]

Nyquist frequency:

\[
\omega=\pi
\]

\end{tcolorbox}
%%%%%%%%%%%%%%%%%%%%%%%%%%%%%%%%%%%%%%%%%%%%%%%%%%%%%%%%%%%%%%%%%%%%%%%%%%%%%%%%
%%%%%%%%%%%%%%%%%%%%%%%%%%%%%%%%%%%%%%%%%%%%%%%%%%%%%%%%%%%%%%%%%%%%%%%%%%%%%%%%
\section{Linearity Property}

If

\[
x_1[n]\xleftrightarrow{\mathcal{F}}X_1(e^{j\omega})
\]

and

\[
x_2[n]\xleftrightarrow{\mathcal{F}}X_2(e^{j\omega}),
\]

then

\[
\boxed{
a_1x_1[n]+a_2x_2[n]
\xleftrightarrow{\mathcal{F}}
a_1X_1(e^{j\omega})+a_2X_2(e^{j\omega})
}
\]

%%%%%%%%%%%%%%%%%%%%%%%%%%%%%%%%%%%%%%%%%%%%%%%%%%%%%%%%%%%%%%%%%%%%%%%%%%%%%%%%

\section{Time-Shifting Property}

A delay of \(n_0\) samples gives

\[
\boxed{
x[n-n_0]
\xleftrightarrow{\mathcal{F}}
e^{-j\omega n_0}X(e^{j\omega})
}
\]

%%%%%%%%%%%%%%%%%%%%%%%%%%%%%%%%%%%%%%%%%%%%%%%%%%%%%%%%%%%%%%%%%%%%%%%%%%%%%%%%

\subsection{Interpretation}

\begin{itemize}
\item Magnitude remains unchanged.
\item Phase changes linearly with frequency.
\end{itemize}

%%%%%%%%%%%%%%%%%%%%%%%%%%%%%%%%%%%%%%%%%%%%%%%%%%%%%%%%%%%%%%%%%%%%%%%%%%%%%%%%

\section{Frequency-Shifting Property}

Multiplication by a complex exponential shifts the spectrum:

\[
\boxed{
x[n]e^{j\omega_0 n}
\xleftrightarrow{\mathcal{F}}
X(e^{j(\omega-\omega_0)})
}
\]

This is the basis of digital modulation.

%%%%%%%%%%%%%%%%%%%%%%%%%%%%%%%%%%%%%%%%%%%%%%%%%%%%%%%%%%%%%%%%%%%%%%%%%%%%%%%%

\section{Time-Reversal Property}

\[
\boxed{
x[-n]
\xleftrightarrow{\mathcal{F}}
X(e^{-j\omega})
}
\]

If \(x[n]\) is real,

\[
X(e^{-j\omega})=X^*(e^{j\omega}).
\]

%%%%%%%%%%%%%%%%%%%%%%%%%%%%%%%%%%%%%%%%%%%%%%%%%%%%%%%%%%%%%%%%%%%%%%%%%%%%%%%%

\section{Conjugation Property}

\[
\boxed{
x^*[n]
\xleftrightarrow{\mathcal{F}}
X^*(e^{-j\omega})
}
\]

For real sequences,

\[
\boxed{
X(e^{-j\omega})=X^*(e^{j\omega})
}
\]

This is called \textbf{Hermitian symmetry}.

%%%%%%%%%%%%%%%%%%%%%%%%%%%%%%%%%%%%%%%%%%%%%%%%%%%%%%%%%%%%%%%%%%%%%%%%%%%%%%%%

\section{Convolution Property}

For

\[
y[n]=x[n]*h[n],
\]

the DTFT is

\[
\boxed{
Y(e^{j\omega})=
X(e^{j\omega})H(e^{j\omega})
}
\]

Convolution in time becomes multiplication in frequency.

%%%%%%%%%%%%%%%%%%%%%%%%%%%%%%%%%%%%%%%%%%%%%%%%%%%%%%%%%%%%%%%%%%%%%%%%%%%%%%%%

\section{Multiplication Property}

If

\[
y[n]=x[n]h[n],
\]

then

\[
\boxed{
Y(e^{j\omega})=
\frac{1}{2\pi}
\int_{-\pi}^{\pi}
X(e^{j\theta})
H(e^{j(\omega-\theta)})
\,d\theta
}
\]

Thus, multiplication in time corresponds to convolution in frequency.

%%%%%%%%%%%%%%%%%%%%%%%%%%%%%%%%%%%%%%%%%%%%%%%%%%%%%%%%%%%%%%%%%%%%%%%%%%%%%%%%

\section{Differentiation in Frequency}

\[
\boxed{
nx[n]
\xleftrightarrow{\mathcal{F}}
j\frac{d}{d\omega}X(e^{j\omega})
}
\]

This property is useful for computing moments of sequences.

%%%%%%%%%%%%%%%%%%%%%%%%%%%%%%%%%%%%%%%%%%%%%%%%%%%%%%%%%%%%%%%%%%%%%%%%%%%%%%%%

\section{Parseval’s Theorem}

For an energy sequence,

\[
E=
\sum_{n=-\infty}^{\infty}|x[n]|^2.
\]

In the frequency domain,

\[
\boxed{
E=
\frac{1}{2\pi}
\int_{-\pi}^{\pi}
|X(e^{j\omega})|^2\,d\omega
}
\]

Energy is conserved across domains.

%%%%%%%%%%%%%%%%%%%%%%%%%%%%%%%%%%%%%%%%%%%%%%%%%%%%%%%%%%%%%%%%%%%%%%%%%%%%%%%%

\section{Energy Spectral Density}

Define

\[
\boxed{
\Psi(\omega)=|X(e^{j\omega})|^2
}
\]

called the \textbf{energy spectral density (ESD)}.

Then

\[
E=
\frac{1}{2\pi}
\int_{-\pi}^{\pi}
\Psi(\omega)\,d\omega.
\]

%%%%%%%%%%%%%%%%%%%%%%%%%%%%%%%%%%%%%%%%%%%%%%%%%%%%%%%%%%%%%%%%%%%%%%%%%%%%%%%%

\section{Example: Energy of \(a^n u[n]\)}

Let

\[
x[n]=a^n u[n],\qquad |a|<1.
\]

%%%%%%%%%%%%%%%%%%%%%%%%%%%%%%%%%%%%%%%%%%%%%%%%%%%%%%%%%%%%%%%%%%%%%%%%%%%%%%%%

\subsection{Time-Domain Energy}

\[
E=
\sum_{n=0}^{\infty}|a|^{2n}
=
\frac{1}{1-|a|^2}.
\]

%%%%%%%%%%%%%%%%%%%%%%%%%%%%%%%%%%%%%%%%%%%%%%%%%%%%%%%%%%%%%%%%%%%%%%%%%%%%%%%%

\subsection{Result}

\[
\boxed{
E=\frac{1}{1-|a|^2}
}
\]

%%%%%%%%%%%%%%%%%%%%%%%%%%%%%%%%%%%%%%%%%%%%%%%%%%%%%%%%%%%%%%%%%%%%%%%%%%%%%%%%

\section{Linear-Phase FIR Systems}

A real FIR filter has linear phase if its impulse response is symmetric:

\[
\boxed{
h[n]=h[N-1-n]
}
\]

or antisymmetric:

\[
\boxed{
h[n]=-h[N-1-n].
}
\]

%%%%%%%%%%%%%%%%%%%%%%%%%%%%%%%%%%%%%%%%%%%%%%%%%%%%%%%%%%%%%%%%%%%%%%%%%%%%%%%%

\section{Symmetric FIR Example}

Let

\[
h[0]=1,\qquad h[1]=2,\qquad h[2]=1.
\]

This sequence is symmetric.

%%%%%%%%%%%%%%%%%%%%%%%%%%%%%%%%%%%%%%%%%%%%%%%%%%%%%%%%%%%%%%%%%%%%%%%%%%%%%%%%

\subsection{Frequency Response}

\[
H(e^{j\omega})=
1+2e^{-j\omega}+e^{-j2\omega}
\]

Factorizing,

\[
H(e^{j\omega})=
e^{-j\omega}
\left(
2+2\cos\omega
\right)
\]

Hence,

\[
\boxed{
\angle H(e^{j\omega})=-\omega
}
\]

which is linear in \(\omega\).

%%%%%%%%%%%%%%%%%%%%%%%%%%%%%%%%%%%%%%%%%%%%%%%%%%%%%%%%%%%%%%%%%%%%%%%%%%%%%%%%

\section{Phase Delay}

\[
\boxed{
\tau_p(\omega)=
-\frac{\angle H(e^{j\omega})}{\omega}
}
\]

For the previous filter,

\[
\tau_p(\omega)=1.
\]

%%%%%%%%%%%%%%%%%%%%%%%%%%%%%%%%%%%%%%%%%%%%%%%%%%%%%%%%%%%%%%%%%%%%%%%%%%%%%%%%

\section{Group Delay}

\[
\boxed{
\tau_g(\omega)=
-\frac{d}{d\omega}\angle H(e^{j\omega})
}
\]

For a linear-phase FIR filter,

\[
\boxed{
\tau_g(\omega)=\frac{N-1}{2}
}
\]

which is constant.

%%%%%%%%%%%%%%%%%%%%%%%%%%%%%%%%%%%%%%%%%%%%%%%%%%%%%%%%%%%%%%%%%%%%%%%%%%%%%%%%

\section{Importance of Constant Group Delay}

Constant group delay means all frequency components are delayed equally.

\begin{itemize}
\item No waveform distortion.
\item Important in audio systems.
\item Important in communication receivers.
\item Important in data transmission.
\end{itemize}

%%%%%%%%%%%%%%%%%%%%%%%%%%%%%%%%%%%%%%%%%%%%%%%%%%%%%%%%%%%%%%%%%%%%%%%%%%%%%%%%

\section{Magnitude Response Classification}

%%%%%%%%%%%%%%%%%%%%%%%%%%%%%%%%%%%%%%%%%%%%%%%%%%%%%%%%%%%%%%%%%%%%%%%%%%%%%%%%

\subsection{Low-Pass}

Large gain near \(\omega=0\).

%%%%%%%%%%%%%%%%%%%%%%%%%%%%%%%%%%%%%%%%%%%%%%%%%%%%%%%%%%%%%%%%%%%%%%%%%%%%%%%%

\subsection{High-Pass}

Large gain near \(\omega=\pi\).

%%%%%%%%%%%%%%%%%%%%%%%%%%%%%%%%%%%%%%%%%%%%%%%%%%%%%%%%%%%%%%%%%%%%%%%%%%%%%%%%

\subsection{Band-Pass}

Large gain over a middle range of frequencies.

%%%%%%%%%%%%%%%%%%%%%%%%%%%%%%%%%%%%%%%%%%%%%%%%%%%%%%%%%%%%%%%%%%%%%%%%%%%%%%%%

\subsection{Band-Stop}

Small gain over a middle range of frequencies.

%%%%%%%%%%%%%%%%%%%%%%%%%%%%%%%%%%%%%%%%%%%%%%%%%%%%%%%%%%%%%%%%%%%%%%%%%%%%%%%%

\section{Solved Example 8.1}

Find the DTFT of

\[
x[n]=\delta[n]+\delta[n-2].
\]

%%%%%%%%%%%%%%%%%%%%%%%%%%%%%%%%%%%%%%%%%%%%%%%%%%%%%%%%%%%%%%%%%%%%%%%%%%%%%%%%

\subsection{Solution}

Using the shift property,

\[
X(e^{j\omega})=
1+e^{-j2\omega}.
\]

Factorizing,

\[
X(e^{j\omega})=
2e^{-j\omega}\cos\omega.
\]

Therefore,

\[
\boxed{
|X(e^{j\omega})|=2|\cos\omega|
}
\]

and

\[
\boxed{
\angle X(e^{j\omega})=-\omega.
}
\]

%%%%%%%%%%%%%%%%%%%%%%%%%%%%%%%%%%%%%%%%%%%%%%%%%%%%%%%%%%%%%%%%%%%%%%%%%%%%%%%%

\section{Solved Example 8.2}

Find the DTFT of

\[
x[n]=\left(\frac12\right)^n u[n].
\]

%%%%%%%%%%%%%%%%%%%%%%%%%%%%%%%%%%%%%%%%%%%%%%%%%%%%%%%%%%%%%%%%%%%%%%%%%%%%%%%%

\subsection{Solution}

\[
X(e^{j\omega})=
\frac{1}{1-\frac12 e^{-j\omega}}.
\]

Magnitude:

\[
\boxed{
|X(e^{j\omega})|=
\frac{1}{\sqrt{1.25-\cos\omega}}
}
\]

Phase:

\[
\boxed{
\angle X(e^{j\omega})=
-\tan^{-1}
\left(
\frac{0.5\sin\omega}{1-0.5\cos\omega}
\right)
}
\]

%%%%%%%%%%%%%%%%%%%%%%%%%%%%%%%%%%%%%%%%%%%%%%%%%%%%%%%%%%%%%%%%%%%%%%%%%%%%%%%%

\section{Solved Example 8.3}

Determine whether

\[
h[n]=\{1,2,1\}
\]

has linear phase.

%%%%%%%%%%%%%%%%%%%%%%%%%%%%%%%%%%%%%%%%%%%%%%%%%%%%%%%%%%%%%%%%%%%%%%%%%%%%%%%%

\subsection{Check Symmetry}

\[
h[0]=h[2]=1.
\]

Therefore,

\[
\boxed{
h[n]\text{ is symmetric}
}
\]

and the filter has linear phase.

%%%%%%%%%%%%%%%%%%%%%%%%%%%%%%%%%%%%%%%%%%%%%%%%%%%%%%%%%%%%%%%%%%%%%%%%%%%%%%%%

\section{Solved Example 8.4}

Find the group delay of the filter

\[
h[n]=\{1,1,1,1\}.
\]

%%%%%%%%%%%%%%%%%%%%%%%%%%%%%%%%%%%%%%%%%%%%%%%%%%%%%%%%%%%%%%%%%%%%%%%%%%%%%%%%

\subsection{Solution}

The filter length is

\[
N=4.
\]

Hence,

\[
\tau_g=
\frac{N-1}{2}
=
\frac{3}{2}.
\]

\[
\boxed{
\tau_g=1.5\text{ samples}
}
\]

%%%%%%%%%%%%%%%%%%%%%%%%%%%%%%%%%%%%%%%%%%%%%%%%%%%%%%%%%%%%%%%%%%%%%%%%%%%%%%%%

\section{Solved Example 8.5}

A discrete-time system has

\[
H(e^{j\omega})=
1-e^{-j\omega}.
\]

Identify the filter type.

%%%%%%%%%%%%%%%%%%%%%%%%%%%%%%%%%%%%%%%%%%%%%%%%%%%%%%%%%%%%%%%%%%%%%%%%%%%%%%%%

\subsection{Magnitude}

\[
|H(e^{j\omega})|=
2\left|\sin\frac{\omega}{2}\right|.
\]

At \(\omega=0\),

\[
|H|=0.
\]

At \(\omega=\pi\),

\[
|H|=2.
\]

Therefore,

\[
\boxed{
\text{High-pass filter}
}
\]

%%%%%%%%%%%%%%%%%%%%%%%%%%%%%%%%%%%%%%%%%%%%%%%%%%%%%%%%%%%%%%%%%%%%%%%%%%%%%%%%

\section{DTFT Property Table}

\begin{table}[H]
\centering
\caption{Important DTFT properties}
\begin{tabular}{|p{5cm}|p{6cm}|}
\hline
\textbf{Time Domain} & \textbf{Frequency Domain}\\
\hline
\(a_1x_1+a_2x_2\) & \(a_1X_1+a_2X_2\)\\
\hline
\(x[n-n_0]\) & \(e^{-j\omega n_0}X\)\\
\hline
\(x[n]e^{j\omega_0 n}\) & \(X(e^{j(\omega-\omega_0)})\)\\
\hline
\(x[-n]\) & \(X(e^{-j\omega})\)\\
\hline
\(x*h\) & \(XH\)\\
\hline
\(xh\) & \(\frac{1}{2\pi}X*H\)\\
\hline
\(nx[n]\) & \(j\dfrac{dX}{d\omega}\)\\
\hline
\end{tabular}
\end{table}

%%%%%%%%%%%%%%%%%%%%%%%%%%%%%%%%%%%%%%%%%%%%%%%%%%%%%%%%%%%%%%%%%%%%%%%%%%%%%%%%

\section{Common GATE Mistakes}

\begin{enumerate}
\item Forgetting the DTFT is periodic with \(2\pi\).
\item Using CTFT formulas directly for discrete-time signals.
\item Missing the factor \(e^{-j\omega n_0}\) in time shifting.
\item Confusing phase delay with group delay.
\item Forgetting symmetry conditions for linear phase.
\item Using \(\omega\) in hertz instead of rad/sample.
\item Forgetting the factor \(1/(2\pi)\) in Parseval’s theorem.
\end{enumerate}

%%%%%%%%%%%%%%%%%%%%%%%%%%%%%%%%%%%%%%%%%%%%%%%%%%%%%%%%%%%%%%%%%%%%%%%%%%%%%%%%

\section{Chapter Summary}

In this chapter, we developed the DTFT and studied its properties, including
linearity, time shifting, frequency shifting, time reversal, convolution and
differentiation.

We introduced Parseval’s theorem and energy spectral density, and then studied
linear-phase FIR systems, phase delay and group delay.

Finally, we solved several GATE-style problems involving DTFT computation,
frequency response and filter classification.

%%%%%%%%%%%%%%%%%%%%%%%%%%%%%%%%%%%%%%%%%%%%%%%%%%%%%%%%%%%%%%%%%%%%%%%%%%%%%%%%

\begin{tcolorbox}[title=One-Minute Revision,colback=blue!5]

\[
X(e^{j\omega})=
\sum_{n=-\infty}^{\infty}x[n]e^{-j\omega n}
\]

\[
x[n]=
\frac{1}{2\pi}
\int_{-\pi}^{\pi}
X(e^{j\omega})e^{j\omega n}d\omega
\]

\[
X(e^{j(\omega+2\pi)})=X(e^{j\omega})
\]

\[
x[n-n_0]\leftrightarrow e^{-j\omega n_0}X(e^{j\omega})
\]

\[
x[n]e^{j\omega_0 n}\leftrightarrow X(e^{j(\omega-\omega_0)})
\]

\[
Y(e^{j\omega})=X(e^{j\omega})H(e^{j\omega})
\]

\[
E=
\frac{1}{2\pi}
\int_{-\pi}^{\pi}|X(e^{j\omega})|^2d\omega
\]

Linear phase:

\[
h[n]=h[N-1-n]
\]

Group delay:

\[
\tau_g=\frac{N-1}{2}
\]

\end{tcolorbox}
%%%%%%%%%%%%%%%%%%%%%%%%%%%%%%%%%%%%%%%%%%%%%%%%%%%%%%%%%%%%%%%%%%%%%%%%%%%%%%%%